\documentclass[11pt]{article}
\usepackage[T1]{fontenc}
\usepackage[utf8]{inputenc}
\usepackage{lmodern,microtype,geometry,amsmath,amssymb,amsthm,mathtools,booktabs}
\usepackage[colorlinks=true,linkcolor=blue,urlcolor=blue,citecolor=blue]{hyperref}
\geometry{margin=1in}
\setlength{\emergencystretch}{3em}
\newtheorem{theorem}{Theorem}
\newtheorem{lemma}[theorem]{Lemma}
\newtheorem{corollary}[theorem]{Corollary}
\newtheorem{proposition}[theorem]{Proposition}
\theoremstyle{remark}\newtheorem{remark}[theorem]{Remark}
\newcommand{\E}{\mathbb E}\newcommand{\Pp}{\mathbb P}
\newcommand{\R}{\mathbb R}\newcommand{\cX}{\mathcal X}
\newcommand{\HH}{\mathsf H}
\hypersetup{pdftitle={Polynomially many total mistakes for Bayesian neural-network learning},pdfauthor={ChatGPT-6 Astra; Claude Fable 5 and Opus 5; Kaarel Hanni; Dmitry Vaintrob}}
\title{Polynomially many total mistakes\\for Bayesian neural-network learning}
\author{ChatGPT-6 Astra \and Claude Fable 5 and Opus 5 \and Kaarel H\"anni \and Dmitry Vaintrob\thanks{Authors are listed in alphabetical order.}}
\date{Research draft, 21 October 2026}
\begin{document}\maketitle
\begin{abstract}
For uniformly sampled Boolean inputs, the strong-prior, above-mean-field,
and mean-field learning results imply a stronger conclusion than eventual
fixed accuracy: the total number of mistakes over an infinite learning
stream has a polynomial bound, independent of learner width. The bound
holds in expectation and with high probability, and applies both to
posterior sampling and to counting a prediction as a failure unless the
correct label has posterior probability greater than $2/3$. The new
statistical argument works for any consistency posterior whose relevant
functions have a short finite description, even if their prior weights are
extremely unequal. Deleting repeated inputs reduces the infinite iid
stream to a uniformly random permutation. A harmonic bound over its random
prefixes then gives an overhead of order the input dimension. We state
precisely the structural facts imported from the neural-network papers,
give explicit budgets, and distinguish this result from an arbitrary-order
mistake guarantee. A prediction-preserving neuron-merging restriction also
retains polynomial dependence on depth. These new cumulative conclusions
have not been Lean verified.
\end{abstract}

\section{Prediction protocol and the result}

Let $\cX$ be a finite set of size $m\ge2$, let
$f:\cX\to\{-1,1\}$, and let $\pi$ be a prior on Boolean functions
with $\pi(f)>0$. Before round $j$, the learner has seen the true labels
of the preceding inputs and forms the exact consistency posterior
$\Pi_{j-1}$. The next input $X_j$ is then presented, a prediction is
made, and its true label $f(X_j)$ is revealed. Every label is revealed,
including after a correct prediction. Define
\[
 q_j=\Pi_{j-1}\{g:g(X_j)\ne f(X_j)\},\qquad
 S=\sum_j q_j.
\]
Thus $S$ is cumulative posterior error probability. If the learner draws
a function from the current posterior at each round, its total number of
incorrect labels is denoted $Z$. The expected values of $Z$ and $S$ agree.
For a fixed $0<\theta<1$, define the confidence-failure count
\[
 F_\theta=\sum_j\mathbf1\{\Pi_{j-1}(g(X_j)=f(X_j))\le\theta\}.
\]
The $>2/3$ confidence convention is $F_{2/3}$. Pointwise,
\begin{equation}\label{eq:confidence}
 F_\theta\le S/(1-\theta).
\end{equation}
Any majority classifier's number of errors is at most $2S$, including
arbitrary decisions at ties. These inequalities allow posterior mixtures;
they require neither a unique minimum-complexity function nor posterior
concentration on one function.

We consider either a uniformly random permutation of $\cX$, used once,
or an infinite iid stream uniform on $\cX$. In the latter protocol every
repeated input has posterior correctness probability one, and duplicating
a hard conditioning constraint leaves the posterior unchanged. The order
of first appearances is a uniformly random permutation. Consequently the
total variables $S,F_\theta,Z$ have the same laws in the two protocols
(using fresh posterior randomness on first appearances). We can therefore
prove all bounds for $m$ rounds without replacement.

For $\cX=\{-1,1\}^d$, the neural-network consequences are:
\begin{itemize}
\item At every fixed depth, the strong-prior learner has a total mistake
bound polynomial in $d$ and bounded-weight target-network complexity,
at every width above that complexity, once its variance is sufficiently
small at that fixed width.
\item At every fixed depth and every fixed $\gamma>1$, the specified
finite-variance learner has the same type of polynomial bound, uniformly
over all admissible widths. Its polynomial may depend on $\gamma$.
\item The finite-variance mean-field tanh learner, $\gamma=1$, also has
the same polynomial bound at every admissible width.
\end{itemize}
High-probability bounds are polynomial in $\log(1/\delta)$ as well.
The theorem does not replace uniform random input order by arbitrary order.

\section{A cumulative bound from a finite relevant class}

Write $\HH_r=\sum_{k=1}^r1/k$. The following theorem is the new statistical
ingredient. In particular, it does not bound mistakes by the negative log
prior probability of the target.

\begin{theorem}[Random-order cumulative errors]\label{thm:finite}
Suppose $A$ is a finite set of Boolean functions containing $f$ and
\begin{equation}\label{eq:relevant}
 \log|A|\le h,\qquad \pi(A^c)\le\beta\pi(f),\qquad 0\le\beta<1.
\end{equation}
Under either uniform input protocol above,
\begin{equation}\label{eq:expectation}
 \E S=\E Z\le B_E:=1+2(h+1)\HH_{m-1}+m\beta.
\end{equation}
For $0<\delta<1$, put
\[
 a=h+\log(2m/\delta),\qquad
 B_\delta=1+2a\HH_{m-1}+m\beta,\qquad
 U_\delta=(e-1)B_\delta+\log(2/\delta).
\]
Then $\Pp\{S>U_\delta\}\le\delta$ and
$\Pp\{Z>U_\delta\}\le\delta$. These are separate probability statements;
replace $\delta$ by $\delta/2$ to require both simultaneously.
Equation~\eqref{eq:confidence} gives the corresponding bounds for every
confidence threshold and for majority errors.
\end{theorem}

\begin{proof}
After any correctly labeled history, the posterior denominator is at least
$\pi(f)$. Hence the posterior mass outside $A$ is at most $\beta$,
uniformly over all histories.

Let $S_n$ now denote the set of the first $n$ inputs of the permutation,
and let
\[
 R_n=\max_{g\in A:\,g=f\text{ on }S_n}
       \frac1m|\{x:g(x)\ne f(x)\}|.
\]
The maximum exists because $f$ is included. A function with $k$ errors
survives a uniform size-$n$ prefix with probability
\[
 \frac{\binom{m-k}{n}}{\binom mn}
 \le(1-k/m)^n\le e^{-nk/m},
\]
with survival probability zero when $n>m-k$. A union bound gives
\begin{equation}\label{eq:prefix-tail}
 \Pp\{R_n>u\}\le\min\{1,e^{h-nu}\},\qquad
 \E R_n\le\frac{h+1}{n}\quad(1\le n<m).
\end{equation}
All errors of a surviving function are among the remaining $m-n$ inputs.
Conditional on the history, the expected posterior error on the next
uniformly chosen remaining input is therefore at most
\begin{equation}\label{eq:conditional-risk}
 \frac m{m-n}R_n+\beta.
\end{equation}
The first round costs at most one. Summing expectations and using
\[
 \sum_{n=1}^{m-1}\frac{m}{n(m-n)}=2\HH_{m-1}
\]
proves~\eqref{eq:expectation}. For posterior sampling, condition additionally
on the current input before taking the expectation of the sampled label.

For high probability, let $E_n=\{R_n\le a/n\}$. By
\eqref{eq:prefix-tail}, the probability that any $E_n$ fails is at most
$(m-1)e^{h-a}\le\delta/2$. Stop counting at the first failed prefix
event. More precisely, replace the round-$n+1$ variable $Y_{n+1}$ by
$Y'_{n+1}=Y_{n+1}\mathbf1\{E_1,\ldots,E_n\}$, with the indicator one
at $n=0$, where $Y$ is either $q$ or the sampled error indicator.
The indicator is known before the next input. The sum of conditional
expectations of $Y'$ is at most $B_\delta$, by
\eqref{eq:conditional-risk}. For any adapted variables $0\le Y'_j\le1$,
the elementary bound
\[
 \E[e^{Y'_j}\mid\mathcal F_{j-1}]
 \le1+(e-1)\E[Y'_j\mid\mathcal F_{j-1}]
 \le\exp((e-1)\E[Y'_j\mid\mathcal F_{j-1}])
\]
shows that
$\exp(\sum_jY'_j-(e-1)\sum_j\E[Y'_j\mid\mathcal F_{j-1}])$
is a nonnegative supermartingale. Markov's inequality bounds
$\Pp\{\sum Y'_j>U_\delta\}$ by $\delta/2$. On all good prefixes,
$\sum Y'_j=\sum Y_j$; a union bound proves the claim.
\end{proof}

For $m=2^d$, take $\beta=1/m$. Then $\HH_{m-1}\le1+d\log2$, so
\begin{equation}\label{eq:polynomial-overhead}
 B_E\le2+2(h+1)(1+d\log2),\qquad
 U_\delta=O((h+d+\log(1/\delta))(d+1)).
\end{equation}
The relevant description bound need only remain polynomial when
$\log(1/\beta)=d\log2$; it need not be independent of the desired tail.

\paragraph{What is uniform in width?}
The numerical bounds below hold for each admissible learner width with
the same constants. The high-probability theorem is applied to a fixed
learner and its posterior randomness. It does not assert that one event
simultaneously controls independently sampled predictions of every learner.

\section{The direct prior-mass bound}

The smaller-width branches of the finite-variance neural-network proofs
use a different, simpler estimate.

\begin{lemma}[Exact Bayesian telescoping]\label{lem:telescoping}
For every input sequence, including adaptively chosen inputs,
\[
 S\le\sum_j-\log(1-q_j)\le-\log\pi(f)=:J.
\]
Thus $F_\theta\le J/(1-\theta)$ deterministically. Posterior sampling
satisfies $\E Z\le J$, and
\[
 \Pp\{Z>(e-1)J+\log(1/\delta)\}\le\delta.
\]
\end{lemma}
\begin{proof}
The product $\prod_j(1-q_j)$ is the prior probability that a function
fits all revealed labels, and is at least $\pi(f)$. Use
$q\le-\log(1-q)$. The sampled-error assertion follows from the same
exponential supermartingale as above, conditional on each revealed query.
Finite-horizon bounds extend to infinite streams by monotone convergence.
\end{proof}

This lemma alone generally costs a factor growing with neural-network
width and inverse prior variance. The relevant-class theorem avoids that
factor at large widths; the lemma covers the remaining widths where its
cost is already polynomial.

\section{Application to the three neural-network regimes}

\subsection{Model and imported structural results}

The domain is $\{-1,1\}^d$. A depth-$L$ network has $L-1$ hidden layers
of width $N$, a constant input coordinate, bounded $1$-Lipschitz hidden
activation, and a linear readout. Define
\[
 \kappa_N(w)=\frac dN\|W^1\|_F^2+
 \sum_{s=2}^{L-1}\|W^s\|_F^2+N\|W^L\|_2^2.
\]
The raw Gaussian is conditioned, before training, on
$|z_w(x)|>1/2$ for every input. It therefore induces a prior on Boolean
sign functions. Training uses exact label agreement. Let $C=C_L(f)$ be
the least width of a depth-$L$ implementation with strict margin $1/2$
and all weights bounded in magnitude by that width.

We use the following proved structural facts from
\cite{strong,above,meanfield}. These are imports, not newly established
Gaussian-compression theorems. Their explicit scalar budgets appear in
Appendix~\ref{app:budgets}.
\begin{enumerate}
\item There is a common finite set containing every sign function whose
minimum weighted energy is at most $J$, with log cardinality at most
$F_L(d,J)$. This is polynomial in $d,J$ at fixed $L$.
Cloning gives $\kappa_N(f)\le B_L(d,C):=d(d+1)C^2+LC^4$.
\item A raw target neighborhood has mass at least
$\exp(-AN^\gamma/2)$ for every $N\ge C$, where
$A=A_L(d,C)\ge1$ is polynomial at fixed depth and the Gaussian variances
are $N^{1-\gamma}/d,N^{-\gamma},N^{-\gamma-1}$.
\item At fixed $\gamma>1$, for $K\ge1$ and widths satisfying
\[
 N^{\gamma-1}\ge8(d+1),\qquad
 N^{\gamma-1}\ge64a_K^{-2},\qquad
 a_K=[8L(\sqrt K+1)^{L-2}]^{-1},
\]
the raw total mass of full-margin functions with minimum energy greater
than $(L+3)K$ is at most $L e^{-KN^\gamma/8}$.
\item For tanh at $\gamma=1$, if $K\ge4$ and $N\ge T_L(d,K)$,
the raw total mass of full-margin functions outside a common finite class
of log cardinality at most $G_L(d,K)$ is at most $L e^{-NK/8}$.
Both $T_L$ and $G_L$ are polynomial at fixed depth.
\end{enumerate}
The last two bounds concern total exceptional mass, not one exceptional
function at a time. This distinction is essential for this application.

\subsection{Strong-prior limit}

In the strong-prior parameterization, the energy coordinates are independent
centered Gaussians of variance $\sigma^2$. At fixed width,
\[
 2\sigma^2\log\mu_{N,\sigma}\{w:g(x)z_w(x)>1/2\ \forall x\}
 \longrightarrow-\kappa_N(g).
\]
There are finitely many Boolean functions. Hence, with
$A_{N,f}=\{g:\kappa_N(g)\le\kappa_N(f)\}$,
\[
 \frac{\pi_{N,\sigma}(A_{N,f}^c)}{\pi_{N,\sigma}(f)}\longrightarrow0.
\]
All functions tied with the target's energy remain in this set. By cloning
and compression, its log cardinality is at most
$h=F_L(d,B_L(d,C))$.

\begin{corollary}[Strong-prior cumulative learning]\label{cor:strong}
At every fixed $L$, for every $N\ge C$ there is $\sigma_0>0$ such that
all $0<\sigma<\sigma_0$ satisfy Theorem~\ref{thm:finite} with
$m=2^d$, $\beta=2^{-d}$, and $h=F_L(d,B_L(d,C))$.
Consequently the total expected and high-probability mistake bounds are
polynomial in $d,C$ (and $\log(1/\delta)$ for high probability), independently
of $N$. The threshold $\sigma_0$ may depend on $N,d,L,f$ and the activation;
it need not be uniform in width. One fixed sufficiently small variance
works for the entire infinite stream.
\end{corollary}

The same bounds hold for any limit of these finite-domain posterior laws
as $\sigma\downarrow0$. No selection of a unique energy minimizer is
required. Ordinary function-MAP prediction also obeys the corresponding
finite-class random-order bound: when $\beta<1$, every function outside
$A_{N,f}$ has smaller prior mass than the consistent target and cannot be
MAP. Ties must be selected independently of the next input. Parameter MAP
and a choice of an unattained minimum-energy parameter are different
notions and are not identified with function MAP here.

\subsection{Every fixed exponent \texorpdfstring{$\gamma>1$}{gamma > 1}}

Set $m=2^d$, $A=A_L(d,C)$, and
\begin{align*}
 K&=4A+8\log(Lm),\qquad a=a_K,\\
 T&=\left\lceil\max\{C,[8(d+1)]^{1/(\gamma-1)},
                         (8/a)^{2/(\gamma-1)}\}\right\rceil,\\
 h&=F_L(d,(L+3)K),\qquad J=AT^\gamma/2.
\end{align*}
At $N\ge T$, the relative exceptional mass is at most
\[
 L\exp\{-(K/8-A/2)N^\gamma\}
 =L(Lm)^{-N^\gamma}\le1/m.
\]
The corresponding finite set necessarily contains $f$, since otherwise
the ratio would be at least one. At $C\le N<T$,
$-\log\pi(f)\le AN^\gamma/2\le J$.

\begin{corollary}[Above-mean-field cumulative learning]\label{cor:above}
For every fixed $L$ and $\gamma>1$, Theorem~\ref{thm:finite} applies
at $N\ge T$ with the preceding $h,\beta=1/m$;
Lemma~\ref{lem:telescoping} applies at $C\le N<T$ with cost $J$.
Taking the maximum of the two budgets gives total expected and
high-probability mistake bounds polynomial in $d,C,\log(1/\delta)$,
uniformly for every $N\ge C$. These are statements at the specified
positive variances. The polynomial's degree may depend on $L,\gamma$;
this proof makes no uniform claim as $\gamma\downarrow1$.
\end{corollary}

\subsection{Mean field}

For tanh at $\gamma=1$, put
\[
 K=4A+8\log(Lm),\qquad T=T_L(d,K),\qquad
 h=G_L(d,K),\qquad J=AT/2.
\]
The same calculation gives relative exceptional mass at most $1/m$ for
$N\ge T$. Below $T$, the direct target cost is at most $J$.

\begin{corollary}[Mean-field cumulative learning]\label{cor:meanfield}
The finite-variance mean-field tanh learner has total expected and
high-probability mistake bounds polynomial in $d,C,\log(1/\delta)$,
uniformly over every $N\ge C$, by combining Theorem~\ref{thm:finite}
and Lemma~\ref{lem:telescoping} with the preceding budgets.
The polynomial degree may depend on depth.
\end{corollary}

\paragraph{Why $2^{-d}$ does not create an exponential sample cost.}
The relative tail is reduced to $2^{-d}$ by increasing the structural
budget by $O(d)$. The log size of the finite relevant class is polynomial
in this budget. One never requests an approximation accuracy of $2^{-d}$
from a compression theorem whose cost is polynomial in inverse accuracy.

\subsection{An extension to ReLU}

The ReLU manuscript~\cite{relu} replaces bounded-activation compression by
product-norm capacity control. Its independently reviewed structural bounds
can be inserted into the same cumulative theorem. For a concrete
depth-$L$, width-$C$ ReLU margin teacher, set
\[
 P_* = \prod_{s=1}^L\|W_*^s\|_F,\qquad
 a_H=1+\sqrt{2(L-1)\log2},\qquad
 v_*=\lceil4e(d+1)P_*^2a_H^2\rceil.
\]
The ReLU strong-prior proof gives a finite relevant class with
$h=v_*\log(m+1)$, independent of learner width. Choose the variance small
enough at each fixed $N\ge C$ that its relative tail is at most $1/m$.
Theorem~\ref{thm:finite} then gives expected total errors $O(v_*d^2)$,
with the stated logarithmic-confidence extension. For a bounded-weight
teacher, $P_*^2\le(d+1)C^{4L-2}$, so this is polynomial at fixed depth.

There is also a specified finite-variance result for ReLU, without a
small-variance limit. The same manuscript proves the following energy
tail and count. Define
\begin{align*}
 K_R&=L(4dP_*^2)^{1/L},\qquad S_R=d+L,\\
 r_R&=\min\{1,[4\sqrt{2L}(\sqrt{K_R/L}+1)^{L-1}]^{-1}\},\\
 B_R&=\tfrac12(\sqrt{K_R}+r_R)^2+
       S_R\log\left(3\max\{1,\sqrt{S_R}/r_R\}\right),\\
 T_R&=4[B_R+(S_R/2)\log2+\log m],\\
 v_R&=\left\lceil\frac{4(d+1)}d(T_R/L)^La_H^2\right\rceil.
\end{align*}
For every $N\ge C$ and every $\gamma\ge2$, the Gaussian with energy
coordinate variance $N^{-\gamma}$, globally margin-conditioned as above,
has relative function mass at most $1/m$ outside a finite class with
$h=v_R\log(m+1)$. Thus Theorem~\ref{thm:finite} gives an all-width total
mistake bound polynomial at fixed depth, uniformly over $\gamma\ge2$.
For one hidden layer, $L=2$, the same statement holds for every
$\gamma\ge1$. The deeper finite-variance range $1\le\gamma<2$ is not
settled by this imported ReLU argument. These assertions use the
ReLU paper's homogeneous balancing and capacity proofs, not an extension
of the bounded-activation compression theorem by assertion.

\section{Depth-polynomial support restrictions}

In this section $C$ denotes the width of a layered \emph{Boolean} circuit
whose stages fit the learner's hidden layers. Constant wires and any
fixed-basis compilation overhead count toward its width and depth. This
is the teacher convention of the neuron-merging results; it is a stronger
requirement than a comparably small arbitrary real-valued neural teacher.

The same statistical proof applies to finite loss classes, not just
Boolean hypotheses. Suppose a posterior of real-output networks admits,
apart from posterior mass at most $\beta$ at every history, a map into
a finite collection of losses $\ell_h:\cX\to[0,1]$ such that
\[
 \ell_h(X_i)=0\text{ on every training input},\qquad
 \mathbf1\{f(x)z_w(x)\le1/4\}\le\ell_h(x)\quad\text{for every }x.
\]
For a loss with mean $u$, at least $mu$ inputs have positive loss, so the
probability that its empirical loss is zero on a size-$n$ random prefix
is at most $e^{-nu}$. Its remaining mean equals $m/(m-n)$ times its
full-domain mean. Every step of Theorem~\ref{thm:finite} therefore applies
unchanged. The map may depend on the retained network and the history;
the finite loss class must be fixed before drawing the input order. Here
the cumulative posterior error variable uses the tested margin event
$fz_w\le1/4$; the resulting bounds also control sign errors.

The neuron-merging papers~\cite{merging,naturality} supply exactly this
form of deterministic compression. Retained networks have a compressed
grid representative whose clipped output differs uniformly by at most
$3/32$. Training uses signed margin $>1/2$, so the representative has
signed clipped output $\ge13/32$ on training inputs. A tested margin error
$fz_w\le1/4$ gives representative signed clipped output at most $11/32$.
Use a ramp loss equal to one at or below $11/32$ and zero at or above
$13/32$. The resulting finite class has log size polynomial jointly in
$d,L,C,c+c^{-1},\log(1/\beta)$, independently of $N$ and the scale
$0<t\le1$. Choosing $\beta=1/m$ proves a total mistake bound polynomial
in depth as well as width for those structural support restrictions.
At fixed remaining parameters, the teacher-width dependence remains
quadratic up to logarithms, with an extra factor polynomial in $d$.

This does not by itself establish a polynomial-depth total mistake bound
under the weaker internal-noise restriction alone. Its current entropy
compression couples predictions with an error probability $\alpha$ and
has size polynomial in $1/\alpha$. Taking $\alpha$ of order $2^{-d}$
would destroy the desired dimension dependence. That is a limitation of
this extension, not a counterexample to such a cumulative theorem. If the
weaker support is additionally globally margin-conditioned, the exact
finite-function mean-field argument still gives the fixed-depth cumulative
result after using its robust target-mass cylinder; it does not preserve
the improved polynomial depth dependence.

\section{Which input assumptions matter?}

\paragraph{Uniform random order is doing real work.}
A class of $m+1$ simple functions can require $m$ mistakes in a prescribed
order despite log cardinality only $\log(m+1)$. Let $f\equiv+1$ and let
$g_i$ differ from $f$ only at point $x_i$. Give $g_i$ prior weight
$R^{m-i+1}$ and $f$ weight one, and reveal $x_1,\ldots,x_m$.
Before $x_i$, $g_i$ has posterior mass greater than $1-1/R$, because the
previous candidates have been eliminated. For $R>3$, every round is a
$2/3$-confidence failure, majority prediction is wrong every time, and
expected posterior-sampling mistakes exceed $(1-1/R)m$.
All functions are individually simple point exceptions on the Boolean
cube. This is a counterexample to an arbitrary-order conclusion from a
finite relevant-class count alone. It is not a counterexample proved for
the particular Gaussian neural-network priors.

An arbitrary iid input distribution is also different from uniform iid.
In this same example, take probabilities proportional to
$1,\rho,\ldots,\rho^{m-1}$. The chance that the first occurrence of
$x_{i+1}$ precedes that of $x_i$ is $\rho/(1+\rho)$. A union bound shows
that the first-appearance order is the displayed adverse order with
probability at least $1-(m-1)\rho$. Thus even iid sampling can be almost
adverse when point probabilities range over sufficiently many scales.
Again, this isolates a limitation of the statistical inference rather than
proving failure of a specific neural-network prior.

\paragraph{A useful nonuniform extension.}
For any iid distribution on a finite support of size $m$ and minimum
positive probability $p_{\min}$, the same relevant-class assumptions give
\[
 \E S\le 3+(h+1)(1+\log T)+m\beta,
 \qquad T=\left\lceil p_{\min}^{-1}\log(2/p_{\min})\right\rceil.
\]
Indeed, expected relevant-class risk after $n$ labels is at most
$(h+1)/n$. It is also at most the expected unseen probability, which is
at most $e^{-np_{\min}}$. Sum the first bound through $T$ and the second
afterwards. The outside-class contribution is at most $\beta$ times the
unseen probability at every history, and the sum of the expected unseen
probabilities equals the support size $m$. In particular, the total bound
stays polynomial whenever $\log(1/p_{\min})$ does. Uniform sampling is a
clean special case, not the only useful one.

\paragraph{What this does not show about a converse.}
Few realized mistakes in one run do not alone imply that the target has a
small circuit: favorable random guesses and incomplete domain coverage
must be excluded. Nor does a general bounded-mistake learner automatically
supply a small circuit unless its predictions have controlled computational
descriptions. The results here establish the upper direction for specified
Bayesian learners; a converse requires an additional argument.

\appendix
\section{Explicit imported neural-network budgets}\label{app:budgets}

For reproducibility, the quantities below are the scalar bounds supplied
by the earlier proofs. Their mathematical roles are stated in Section~4.
They are intentionally loose.

\paragraph{A common norm-sublevel count.}
For $J\ge1$, put
\begin{align*}
 \rho_J&=\frac{2^{-L}(J+1)^{-(L-1)/2}}{12(L-1)},\\
 u&=\max\left\{\lceil9(J+1)\rceil,
       \left\lceil256\rho_J^{-4}[d+\log(L/\rho_J)]\right\rceil\right\},\\
 F_L(d,J)&=[u(d+2)+(L-2)u^2]\,
 \log\left(3+16Lu\max\{u,d+1\}[u(u+1)]^{L-1}\right).
\end{align*}
The norm-compression theorem preserves all Boolean signs and yields a
width-$u$, weight-bounded implementation. Uniform parameter rounding gives
the count $F_L$. The class is common to all original learner widths.

\paragraph{The target cylinder.}
Set
\begin{align*}
 b_L&=\sqrt{16\log9+8\log(4L)},\qquad a_0=2+3C^2+b_L,\\
 \rho_0&=[2+9C^2(d+L+1)a_0^{L-2}]^{-1},\\
 B_0&=d(d+1)C^2+2(L-2)C^4+36C^4,\\
 A_L(d,C)&=2B_0+4[d+2+(L-2)C]/\rho_0+2.
\end{align*}
The earlier cylinder proof gives raw target mass at least
$\exp(-A_L(d,C)N^\gamma/2)$ for every $\gamma\ge1,N\ge C$.

\paragraph{Mean-field exact-function compression.}
Let $r=L-1$, $K\ge4$, and choose backwards
\[
 g_r=32rK^2,\qquad
 g_s=8g_{s+1}(1+g_{s+1})^2\quad(s=r-1,\ldots,1).
\]
Put
\[
 q=\lceil g_1^2/2\rceil,\quad
 \ell=\log(2\cdot2^{2d}+1),\quad
 T_L(d,K)=\left\lceil\max\{1,2q\ell/(K-1)\}\right\rceil.
\]
For the finite-description bound, write $g=g_1$, $b=8(1+g)^2$, and set
\begin{align*}
 \eta&=[32Kr b^{r-1}]^{-1},&
 u_*&=\lceil2\eta^{-2}\log(8r\,2^{2d})\rceil,\\
 D_*&=u_*(d+1)+(L-2)u_*^2,&
 B_*&=\max\{1,\sqrt{4g+2\log(8\max\{D_*,1\})},4K/u_*\},\\
 P_*&=D_*+u_*,&
 \Lambda_*&=(d+Lu_*)[1+u_*(B_*+1)]^{L-1}.
\end{align*}
Then $G_L(d,K)=P_*\log(3+16B_*\Lambda_*)$ is valid. At fixed depth,
$T_L$ and $G_L$ have polynomial majorants in $d,K$. The covariance
compression uses the global margin to recover exactly the whole Boolean
sign function. It is not merely an average-error approximation.

\begin{thebibliography}{9}
\bibitem{strong} ChatGPT-6 Astra, Claude Fable 5 and Opus 5, Kaarel H\"anni,
and Dmitry Vaintrob. \emph{Width-independent generalization in the
strong-prior limit of Bayesian neural networks}.
Research manuscript, 2026.
\url{https://github.com/kaarelh/AI-works/tree/main/neural-network-generalization/main-results/strong-prior-generalization}.
\bibitem{above} ChatGPT-6 Astra, Claude Fable 5 and Opus 5, Kaarel H\"anni,
and Dmitry Vaintrob. \emph{Width-uniform Bayesian generalization above the
mean-field scaling}. Research manuscript, 2026.
\url{https://github.com/kaarelh/AI-works/tree/main/neural-network-generalization/main-results/gamma-above-one}.
\bibitem{meanfield} ChatGPT-6 Astra, Kaarel H\"anni, and Dmitry Vaintrob.
\emph{Mean-field Bayesian circuit learning at every width above target
complexity}. Research manuscript, 2026.
\url{https://github.com/kaarelh/AI-works/tree/main/neural-network-generalization/main-results/meanfield-all-width}.
\bibitem{merging} ChatGPT-6 Astra, Claude Fable 5 and Opus 5, Kaarel H\"anni,
and Dmitry Vaintrob. \emph{Quadratic-width learning bounds for Bayesian
networks with stable neuron populations}. Research manuscript, 2026.
\url{https://github.com/kaarelh/AI-works/tree/main/neural-network-generalization/works-in-progress/improved-polynomials/neuron-averaging}.
\bibitem{naturality} ChatGPT-6 Astra, Claude Fable 5 and Opus 5, Kaarel H\"anni,
and Dmitry Vaintrob. \emph{Quadratic-width learning with symmetry-invariant
neuron merging}. Research manuscript, 2026.
\url{https://github.com/kaarelh/AI-works/tree/main/neural-network-generalization/works-in-progress/improved-polynomials/neuron-merging-naturality}.
\bibitem{relu} ChatGPT-6 Astra, Claude Fable 5 and Opus 5, Kaarel H\"anni,
and Dmitry Vaintrob. \emph{Polynomial sample complexity for ReLU NNBayes
with strong Gaussian priors}. Research manuscript, 2026.
\end{thebibliography}
\end{document}
